\section{Related Work in RAG Groundedness}
\label{sec:related_work}

The three pillars of RAG are: (1) context relevance (2) groundedness and (3) answer relevance. For this work, we focus on the second pillar, groundedness, which is the degree to which the model's output is supported by the retrieved context. Groundedness is a key component of pluralistic reasoning in RAG pipelines, as it ensures that the model's output is based on a diverse set of information sources. Work on this topic splits along the epistemic--aleatoric axis, with epistemic work focusing on factuality type questions where there is a single ground truth answer, and aleatoric work focusing on opinion type questions where there are multiple valid answers.

\paragraph{Epistemic Groundedness.} Work done by \citet{ravichander-etal-2022-condaqa} investigates whether a model's output is supported by the retrieved context under paraphrase, scope, and negation perturbations. They find that models are not robust to these perturbations, and are particularly sensitive to negation and changes to scope. \citet{qi-etal-2024-long2rag} investigate the robustness of RAG pipelines, and find that models are not robust to long context, and that the performance is varies greatly across domain (e.g. History, AI, Biology) and query type (e.g. factual, methodological, explanatory).

\paragraph{Aleatoric Groundedness.} \citet{wan-etal-2024-evidence} investigate the convincingness of evidence in RAG pipelines (i.e. whether a piece of selected evidence is used by the model) for contentious queries (e.g. Are humans fundamentally good or evil?), and find that language models do not always align with human values when evaluating evidence. They analyse convincingness by asking a language model a yes or no question and presenting it with a single piece of evidence for each position. As shown by \citet{naphade-2026-rational}, convincingness additionally changes when multiple pieces of evidence are presented for each position. \citeposs{wu2024clasheval} work further complicates this issue by illustrating how even frontier models ignore RAG evidence when it is presented in a way that is not aligned with the model's prior beliefs. 

We see that across the epistemic--aleatoric axis, RAG pipelines are not robust to perturbations in the retrieved context, and that this lack of robustness is model-agnostic, domain and query type dependent, and is particularly pronounced for aleatoric queries. 